\frontmatter
\begin{titlepage}
\thispagestyle{empty}
\begin{tikzpicture}[remember picture,overlay]
\node at (current page.center) {\origincludegraphics[width=\paperwidth,height=\paperheight]{art/cover.png}};
\node[anchor=north,text width=0.86\paperwidth,align=center] at ([yshift=-2.6cm]current page.north)
  {\color{white}\sffamily\bfseries\fontsize{42}{50}\selectfont Potentials,\hspace{0.35em}Primes\\[4pt] and\hspace{0.35em}Prediction};
\node[anchor=north,text width=0.80\paperwidth,align=center] at ([yshift=-6.0cm]current page.north)
  {\color{white!85!cyan}\sffamily\fontsize{15}{20}\selectfont A First Course on the Riemann Zeros,\\ Quantum Spectra and Learning from Data};
\node[anchor=south,text width=0.80\paperwidth,align=center] at ([yshift=1.9cm]current page.south)
  {\color{white}\sffamily\fontsize{14}{18}\selectfont Directed by Dharmaj Soni\\[2pt]{\color{white!70}\fontsize{11}{14}\selectfont Written with Claude (Anthropic)}};
\end{tikzpicture}
\end{titlepage}
\chapter*{About this book}
This book takes a reader who has completed a first-year university physics course to the point where they can read, line by line, the research paper ``Fitted potentials on the Wu--Sprung background memorize the Riemann zeros while the explicit formula predicts them'' \cite{soniR1}, and rerun its key computations. Along the way it is a self-contained course in complex numbers, Fourier analysis, linear algebra, one-dimensional quantum mechanics, the Riemann zeta function, random matrices, numerical methods, and the statistics of prediction.

\section*{How this book was made}
The book was planned and directed by Dharmaj Soni, the author of the paper. Its chapters were drafted by the large language models Claude Opus 5.5 and reviewed by Claude Fable 5.2 (Anthropic), following a written specification (the ``book bible'', included in the source repository). Every numerical statement in a worked example or answer is recomputed by a verification script that is part of the source, and every figure is generated by a script in the repository. AI-written text can contain errors that such checks do not catch; corrections are welcome through the repository.

\section*{Pictures and facts}
Every chapter opens with a picture computed from that chapter's own mathematics (the scripts are in the repository folder \texttt{tex/figs/art}); none is a stock image. The ``Did you know?'' boxes, roughly one every ten pages, were drafted by Claude and checked by the coordinating editor; the source behind each fact is recorded in \texttt{bible/did\_you\_know.json}.

\section*{Navigating the PDF}
Every entry in the Contents, every cross-reference and every citation is a link. Each chapter's opening page has a link back to the Contents, and the PDF bookmarks panel lists all parts, chapters and sections.
